% SOURCE AUTHORITY:
% expose_VII_body.tex lines 1010--1110; scan indices 190--194; source folios 179--183.
% Direct scan comparison: physical PDF pages 191--195 at 600 dpi; diagram 3.1.3 checked
% node-by-node and arrow-by-arrow against physical page 192.
% Transparent notation corrections: printed/typewritten subscript letter o is rendered as zero
% throughout L_0, d_0, H_0 and degree 0; the erroneous prose reference 3.1.7.1 is 3.1.8.
% Hard stop before source line 1116, section 3.2.

\section*{3. Supplements on homological algebra}

The principal purpose of this section is to define the isomorphism
\[
 \Biext^1(P,Q;G)\simeq
 \Ext^1(P\overset{\mathbb L}{\otimes}Q,G)
\]
announced above. For this we shall need certain developments in
homological algebra which are of independent interest and will occupy
sections 3.1 through 3.4. The developments in 3.1 and 3.3 may be
regarded as supplements to [4], especially to section 6.9 therein.

\subsection*{3.1. Construction of an additive cofibred category
$\underline\Psi_{L_\bullet}$ in terms of a chain complex $L_\bullet$
in an abelian category $\underline C$}

Let $\underline C$ be an abelian category and let $L_\bullet$ be a
chain complex in $\underline C$,
\begin{equation}\tag{3.1.1}
 L_2\longrightarrow L_1\longrightarrow L_0\longrightarrow0,
\end{equation}
written with lower-degree conventions, so that the differential has
degree $-1$. Thus $L_i=0$ for $i<0$. We shall define a cofibred
category
\[
 \underline\Psi=\underline\Psi_{L_\bullet}
\]
over $\underline C$; up to canonical equivalence of cofibred
categories it is the category of the same name in [4, 6.9]. The
objects of $\underline\Psi_{L_\bullet}$ are pairs
\[
 (X,\alpha),
\]
where $X$ is an extension of $L_0$ by an object $A$ of
$\underline C$, depending on $X$,
\begin{equation}\tag{3.1.2}
 0\longrightarrow A\xrightarrow{i}X\xrightarrow{j}L_0
 \longrightarrow0,
\end{equation}
and $\alpha$ is a trivialization of the inverse image $d_0^*(X)$ of
this extension by $d_0:L_1\longrightarrow L_0$. It is required that
the corresponding trivialization of the inverse image
$d_1^*d_0^*(X)$ by $d_1:L_2\longrightarrow L_1$ be the evident one,
obtained from the transitivity isomorphism
\[
 d_1^*d_0^*(X)\simeq(d_0d_1)^*(X)
\]
and the relation $d_0d_1=0$.

Equivalently, $\alpha$ is a lifting of $d_0:L_1\longrightarrow L_0$
to a morphism $\alpha:L_1\longrightarrow X$, subject to
$\alpha d_1=0$:
\begin{equation}\tag{3.1.3}
\begin{tikzcd}[row sep=large,column sep=large]
 &&0\arrow[d]\\
 &&A\arrow[d,"i"]\\
 &&X\arrow[d,"j"]\\
 L_2\arrow[r,"d_1"']&L_1\arrow[r,"d_0"']\arrow[ru,"\alpha"]
   &L_0\arrow[d]\\
 &&0
\end{tikzcd}
\qquad,\qquad \alpha d_1=0.\quad(*)
\end{equation}
Here is a simpler description, pointed out by Deligne:
$\underline\Psi_{L_\bullet}(A)$ is the category of extensions of the
complex $L_\bullet$ by $A$, regarded as a complex concentrated in
degree zero.

In this form one sees that the description of the objects
$(X,\alpha)$ depends only on the chain complex $[L_\bullet]$
\underline{truncated in degree 1}, obtained from $L_\bullet$ by
``killing $H_i(L_\bullet)$ for $i\geq2$'': its terms are zero in
degrees at least $2$, agree with those of $L_\bullet$ in degrees at
most $0$, and its degree-$1$ term is the cokernel of
$d_1:L_2\longrightarrow L_1$.

A homomorphism from $(X,\alpha)$ to $(X',\alpha')$ is a homomorphism
of extensions $u:X\longrightarrow X'$, compatible with the identity
on $L_0$, such that $u\alpha=\alpha'$. Homomorphisms are composed by
composing the corresponding homomorphisms of extensions. This gives
the announced category $\underline\Psi_{L_\bullet}$. Up to canonical
isomorphism, it depends only on the degree-$1$ truncation
$[L_\bullet]$ of $L_\bullet$.

There is an evident functor $(X,\alpha)\longmapsto A$,
\begin{equation}\tag{3.1.4}
 \pi:\underline\Psi_{L_\bullet}\longrightarrow\underline C.
\end{equation}
\underline{This functor is cofibred}. Indeed, let
$u:A\longrightarrow B$ be a morphism of $\underline C$, and let
$(X,\alpha)$ be an object of the fibre
$\underline\Psi_{L_\bullet}(A)$. Thus $X$ is an extension of $L_0$ by
$A$. Form the extension $u_*(X)$ of $L_0$ by $B$, the direct image of
$X$ under $u$, and equip it with the partial trivialization
$\beta=f\alpha$, where $f:X\longrightarrow u_*(X)$ is the canonical
morphism. It follows immediately that
\[
 f:(X,\alpha)\longrightarrow(u_*(X),\beta)
\]
is cocartesian for (3.1.4): it has the usual universal property for
$u$-morphisms in $\underline\Psi_{L_\bullet}$ with source
$(X,\alpha)$. This is just the familiar and elementary corresponding
universal property of $f:X\longrightarrow u_*(X)$.

Finally, the cofibred category $\underline\Psi_{L_\bullet}$ is
\underline{additive}; that is, it satisfies
\begin{equation}\tag{3.1.5}
 \underline\Psi_{L_\bullet}(0)
 \quad\text{is equivalent to the one-point category},
\end{equation}
and, for arbitrary $A,B$ in $\underline C$,
\begin{equation}\tag{3.1.6}
 \underline\Psi_{L_\bullet}(A\times B)
 \longrightarrow
 \underline\Psi_{L_\bullet}(A)\times
 \underline\Psi_{L_\bullet}(B)
 \quad\text{is an equivalence}.
\end{equation}
This follows immediately from the corresponding properties of the
cofibred category of extensions of $L_0$ by variable objects of
$\underline C$. We shall freely use the consequences of the
additivity of $\underline\Psi_{L_\bullet}$ established in [4, 1] and
recalled above in 2.5.

3.1.7. Among these consequences, for every object $A$ of
$\underline C$ the category $\underline\Psi_{L_\bullet}(A)$ has a
\underline{zero}, or \underline{trivial}, object (2.5.1). It is the
trivial extension $X=L_0+A$ of $L_0$ by $A$, equipped with the
partial trivialization $\alpha:L_1\longrightarrow X$ whose components
are $d_0$ and $0$.

Moreover, by 2.5.2, for every $(X,\alpha)$ in
$\underline\Psi_{L_\bullet}(A)$ the automorphism group of this object
is canonically isomorphic to the automorphism group of the trivial
object. Denote the latter commutative group by
$\Psi^0_{L_\bullet}(A)$. The explicit description above gives the
canonical isomorphisms
\begin{multline}\tag{3.1.8}
 \Psi^0_{L_\bullet}(A)
 =\Ker\bigl(\Hom(L_0,A)\longrightarrow\Hom(L_1,A)\bigr)
 \simeq\Hom(H_0(L_\bullet),A)\\
 \simeq\Ext^0(L_\bullet,A),
\end{multline}
induced by the usual identification of the automorphism group of the
extension $X$ with $\Hom(L_0,A)$. The isomorphism (3.1.8) is plainly
functorial in $A$.

3.1.9. Finally, recall from 2.5.2 that the natural multiplicative
structure on the fibre category $\underline\Psi_{L_\bullet}(A)$
endows its set of isomorphism classes with a natural group structure.
We denote this group by $\Psi^1_{L_\bullet}(A)$; it is an additive
functor of $A$.

